% Written for the 2026-11-12 ConciseIPRA investigation.
% A self-contained mathematical appendix; no machine formalization is claimed.
% Requires amsmath, amssymb, amsthm and theorem/lemma/proof environments.

\section{The progress rule cannot simply be deleted}
\label{sec:progress-independence}

Let $\mathcal N^{-}$ be the native scalar kernel with Hyp, the two
iterator computation equations, equality transport, QE, DE and DI by
cases, but without QI by progress.  Its raw terms and formulas are
\[
 t::=x\mid0\mid S(t)\mid I(n,a;u.h),\qquad
 A::=(s=t)\mid Q_y A\mid D_{x,y}A.
\]
The intended readings remain
\[
 Q_y B=B(0)\to\forall y B(y),\qquad
 D_{x,y}B=\exists x\bigl(B(x,0)\to\forall y B(x,y)\bigr).
\]
We use an ordinary classical structure to refute derivability.  A
classical structure is also a one-world intuitionistic structure, so
this is a legitimate independence argument for the intuitionistic
calculus.

\begin{theorem}[Independence of equation progress]
\label{thm:progress-independent}
Put $E(y)\equiv I(y,0;u.u)=0$.  The full native kernel proves
$Q_yE(y)$, $E(v)$ for a free numerical variable $v$, and the closed
public Eq/$D$ formula
\[
                 D_{x,y}\bigl(I(y,0;u.u)=0\bigr),
                 \qquad x\text{ unused}.
\]
The kernel $\mathcal N^{-}$ proves none of them.  This remains true even if the
ordinary target theory of $\mathcal N^{-}$ is augmented by successor
injectivity, coverage, $0\ne S0$, absence of all finite successor
cycles, and extensionality of fixed iterator bodies.
\end{theorem}

The construction below adapts the saturation and elementary-chain
method used by Friedman for primitive recursion without induction
\cite[Theorem~8 and Lemmas~24--25]{friedman1975}.  The extra coherence
condition is stated explicitly because the native syntax uses a bound
body, rather than an independent name for every recursion program.
Consequently this appendix claims an application to the native kernel,
not a new general model-theoretic method.

\begin{lemma}[An orbit in a saturated structure]
\label{lem:biinfinite-orbit}
If $M$ is countably saturated and $H:M\to M$ is a total function
definable with finitely many parameters, there is a sequence
$(a_j)_{j\in\mathbb Z}$ in $M$ satisfying $H(a_j)=a_{j+1}$.
\end{lemma}
\begin{proof}
The partial type
\[
 \{\exists z\ H^m(z)=x:m\in\mathbb N\}
\]
is finitely satisfiable: for a finite collection with maximum exponent
$m$, choose $x=H^m(z_0)$ for any $z_0$.  Saturation supplies $a_0$ in
all finite iterate images.  Given such an $a_{-k}$, the type in $x$
consisting of $H(x)=a_{-k}$ and the same image conditions is finitely
satisfiable, since $a_{-k}$ lies in the image of $H^{m+1}$ for every
$m$.  Choose $a_{-k-1}$ realizing it.  Define the positive-index
elements by forward iteration.  Only countable types in one variable
and finite parameter sets have been used.
\end{proof}

\paragraph{The coherence requirement.}
Use a first-order description language.  For a finite update term
$h(u,\vec p)$ in earlier symbols, its iterator has a symbol
$J_h(n,a,\vec p)$, with equations
\begin{align}
 J_h(0,a,\vec p)&=a,\label{eq:independence-zero}\\
 J_h(Sn,a,\vec p)&=h(J_h(n,a,\vec p),\vec p).
 \label{eq:independence-step}
\end{align}
For every two update templates whose iterator symbols have been
introduced, require
\begin{multline}
 \forall\vec p\,\forall\vec q\,
 \bigl(\forall u\ h(u,\vec p)=k(u,\vec q)\bigr)\ \longrightarrow\\
 \forall n\,\forall a\ J_h(n,a,\vec p)=J_k(n,a,\vec q).
 \label{eq:independence-coherence}
\end{multline}
This is an external family of first-order sentences, not an
object-language function quantifier.  It is stronger than the
substitution coherence needed below.

\begin{lemma}[Coherent extension by iterators]
\label{lem:coherent-iterator-extension}
Suppose $M$ is a countably saturated structure with an injective
successor, $0$ outside the successor range, coverage, and no finite
successor cycles.  Suppose an existing family of fixed-body iterator
symbols satisfies \eqref{eq:independence-zero}--\eqref{eq:independence-coherence}.
The same domain can be expanded by iterator symbols for every update
term over its present language, preserving all three requirements and
all previously interpreted symbols.
\end{lemma}
\begin{proof}
The successor graph is one component of type $\mathbb N$, containing
$0$, and components of type $\mathbb Z$.  Choose an origin in each
$\mathbb Z$ component.

For each actual unary function $H$ defined by a present update term
and a tuple of parameters, first ask whether $H$ is represented by an
already interpreted iterator template at some parameter tuple.  If it
is, retain that iterator's values for every counter and seed.  Condition
\eqref{eq:independence-coherence} makes this independent of the chosen
old representation.

For each remaining $H$, on the component of $0$ use the standard
finite iterations $H^m(a)$ starting at the supplied seed $a$.  On
each $\mathbb Z$ component choose a bi-infinite $H$-orbit using
Lemma~\ref{lem:biinfinite-orbit}, and assign its entries to the counter
component in successor order.  The orbit may be chosen separately
for each seed; choosing the same orbit for all seeds is also allowed.

The choices are indexed by actual functions $H$, not by their
syntactic descriptions.  Thus two new parameterized templates
representing the same $H$ receive the same iterator values, and a
new template equal to an old one inherits the old values.  This
proves coherence, including new--old pairs.  The zero equation holds
in the $\mathbb N$ component; the successor equation holds there by
finite iteration and in every other component by the orbit equation.
Every required update function is definable in the structure before
the expansion, so no circular interpretation has been used.
\end{proof}

\paragraph{An initial failure of induction.}
Let $M_0$ have underlying set the disjoint union
$\mathbb N\sqcup\mathbb Z$.  Successor is ordinary successor on
each component, and zero belongs to $\mathbb N$.  Introduce the
identity-update iterator $J(n,a)$ by
\[
 J(n,a)=\begin{cases}
 a,&n\text{ belongs to the }\mathbb N\text{ component},\\
 Sa,&n\text{ belongs to the }\mathbb Z\text{ component}.
 \end{cases}
\]
Then $J(0,a)=a$ and $J(Sn,a)=J(n,a)$, exactly the computation
equations for $h(u)=u$.  There is only one iterator template, so
coherence is automatic.  For an element $c$ of the $\mathbb Z$
component, $J(c,0)=S0\ne0$.

\paragraph{Closing the language.}
Starting with $M_0$, alternate the following two operations:

\begin{enumerate}
\item Take a countably saturated elementary extension in the current
first-order language.
\item Apply Lemma~\ref{lem:coherent-iterator-extension} to add all
missing iterator templates over that language.
\end{enumerate}

Use inclusions for the elementary embeddings and take the union after
countably many stages.  Each language is countable, since its
descriptions are finite strings.  Every raw scalar term is included
at some finite stage, by induction on its finite syntax.  Each fixed
computation or coherence sentence holds at its stage of introduction,
is preserved at all later elementary extensions, and therefore holds
in the union.  More formally, for a fixed stage language the later
reducts form an elementary chain.  Expanding a structure does not
change the values of old symbols.  The equation $J(c,0)=S0$ and
$S0\ne0$ are preserved throughout.

Saturation may create new counter components on which $J(n,0)=0$;
the construction does not assert that this equation defines exactly
the externally standard numbers in the final model.  Only the
preserved counterexample $c$ is needed.

\paragraph{Why this interprets the raw binder syntax.}
Translate a raw iterator $I(n,a;u.h)$ using the symbol for the finite
translated body $h$ and its passive parameters.  This translation
respects capture-avoiding substitution semantically.  At the iterator
case, the two ways of translating a substitution may use different
template symbols, but the structural induction hypothesis says their
actual unary update functions are pointwise equal.  Condition
\eqref{eq:independence-coherence} equates the two iterator values.
The variable, zero and successor cases are immediate.  It follows
that equality transport in arbitrary raw term or formula templates
is sound, including templates with nested iterator binders.

\begin{proof}[Proof of Theorem~\ref{thm:progress-independent}]
In the constructed structure read native equality as actual equality
and read $Q,D$ by their displayed ordinary formulas.  The computation
rules hold by construction.  Hyp and transport are sound by ordinary
equality and substitution.  QE and DE are sound by their ordinary
logical readings.  DI by cases is sound because every numerical
element is zero or a successor; its eigenvariable condition permits
universal generalization in the successor branch.  Thus all rules of
$\mathcal N^{-}$ are sound.

The zero instance $E(0)$ is true.  At the preserved element $c$,
$I(c,0;u.u)=S0\ne0$, so $\forall y E(y)$ and $Q_yE(y)$ are false.
The open formula $E(v)$ is false under the assignment $v=c$.
For an unused witness variable $x$, the ordinary reading of
$D_{x,y}E(y)$ is $\exists x\,Q_yE(y)$.  The numerical domain is
whole and nonempty, with the designated element $0$, so this is
equivalent to $Q_yE(y)$ and is false as well.  None of the three
formulas can therefore be derived in $\mathcal N^{-}$.

In the full kernel, zero computation proves $E(0)$, and successor
computation gives
\[
 I(Sz,0;u.u)=I(z,0;u.u).
\]
Transport under the assumption $E(z)$ gives $E(Sz)$.  QI by progress
now proves $Q_yE(y)$, and QE with $E(0)$ proves $E(v)$.
In this QI profile, $D_{x,y}E(y)$ is also obtained by DI by cases
at witness $0$: under the temporary assumption $E(0)$, QE of the
proved closure supplies $E(Sz)$.  Equivalently, this is the usual
native equivalence between $Q_yE(y)$ and its unused-witness $D$
form.  In the original Choice-DI profile, whose progress rule
concludes a $D$ formula, the same progressive branch directly
introduces $D_{x,y}E(y)$ at witness $0$.  Thus the separator is in
the public Eq/$D$ language and does not depend on which of these
equivalent progress-rule profiles is used.
\end{proof}

\paragraph{Scope of the obstruction.}
The result concerns deleting the progress principle while retaining
the other stated primitives and their ordinary readings.  It is not
a lower bound on the size of an arbitrary bi-interpretable
presentation.  In particular it does not obstruct fusing progress
with another rule, changing the term language, or transporting
induction through a different interpretation.  The construction is
a written model-theoretic proof; it has not been checked by Lean or
Rocq.

% Bibliography entry for the including document:
% \bibitem{friedman1975}
% Harvey Friedman, Provable equality in primitive recursive arithmetic
% with and without induction, Pacific Journal of Mathematics 57(2)
% (1975), 379--392. DOI: 10.2140/pjm.1975.57.379.
% https://msp.org/pjm/1975/57-2/pjm-v57-n2-p10-s.pdf
